\subsection{Coprime elements}

DEFINITION 5. --- In an ordered group two elements x and y are called coprime if \(\inf(x, y) = 0\) .

In some cases it is natural to define two elements to be coprime if \(\inf(|x|, |y|) = 0\) (cf.~INT, II, § 1) or to introduce the corresponding terminology in divisibility theory. We shall not do so here.

Two coprime elements are necessarily positive. The positive and negative parts \(x^{+}\) and \(x^{-}\) of x are coprime (VI, p.~12, Prop. 9, a)). The elements \(x_{i}\) of a family \((x_{t})_{t\in I}\) are said to be setwise coprime if \(\inf_{t\in I}x_t = 0\) ; if the \(x_{t}\) are all \(\geqslant 0\) then it is

sufficient for there to exist a finite subset J of I such that the corresponding elements are setwise coprime. The elements of a family \((x_{\iota})\) are said to be pairwise coprime if \(\inf(x_{\iota}, x_{\kappa}) = 0\) for every pair \((\iota, \kappa)\) of distinct indices.

The \(x_{i}\) can be setwise coprime without being pairwise coprime.

If \(x\) and \(y\) are coprime, we also say that \(x\) is coprime to \(y\) , or that \(y\) is coprime to \(x\) .

(DIV) Two elements x and y of K are said to be coprime if the principal ideals \((x)\) and \((y)\) are nonzero and coprime in \(P^{*}\) ; this amounts to saying that 1 is a gcd of x and y, and implies that x and y belong to A. For example the numerator and denominator of a reduced fraction are coprime. The notions of pairwise and setwise coprime elements are defined similarly.

(DIV) When \(x\) and \(y\) are coprime, they are often said to be « prime to one another »; it is convenient to avoid this terminology, which can lead to confusion with the notion of prime numbers (I, p.~50, Def. 16).

PROPOSITION 10. --- Let \(x, y\) and \(z\) be three elements of an ordered group; for \(x - z\) and \(y - z\) to be coprime, it is necessary and sufficient that \(z = \inf(x, y)\) .

Indeed the relations \(z = \inf (x, y)\) and \(0 = \inf (x - z, y - z)\) are equivalent.

PROPOSITION 10 (DIV). --- Let \(a, b\) and \(c\) be three elements of \(\mathbf{K}\) with \(c \neq 0\) ; for the quotients \(ac^{-1}\) and \(bc^{-1}\) to be coprime, it is necessary and sufficient that \(c\) be a gcd of \(a\) and \(b\) .

PROPOSITION 11. --- If \((x_{i})\) , \((y_{j})\) are two finite families of positive elements of a lattice ordered group, then

\[
\inf \left(\sum_ {i} x _ {i}, \sum_ {j} y _ {j}\right) \leqslant \sum_ {i, j} \inf (x _ {i}, y _ {j}).
\]

Arguing by induction on the number of elements in the families \((x_{i})\) and \((y_{j})\) , it is enough to prove that, if \(x, y\) and \(z\) are positive elements, then

\[
\inf (x, y + z) \leqslant \inf (x, y) + \inf (x, z).
\]

Indeed, put \(t = \inf(x, y + z)\) ; by VI, p.~12, Cor., we can write \(t = t_1 + t_2\) with \(0 \leqslant t_1 \leqslant y\) and \(0 \leqslant t_2 \leqslant z\) ; since \(t_1\) and \(t_2\) are positive, we also have \(t_1 \leqslant x\) and \(t_2 \leqslant x\) , whence \(t_1 \leqslant \inf(x, y)\) and \(t_2 \leqslant \inf(x, z)\) .

COROLLARY 1. --- If x and y are two coprime elements, and z a positive element, of a lattice ordered group, then \(\inf(x, z) = \inf(x, y + z)\) .

Indeed \(\inf (x,y + z)\leqslant \inf (x,z)\) by Prop. 11, and \(\inf (x,z)\leqslant \inf (x,y + z)\) since \(y\geqslant 0\) , whence the corollary.

COROLLARY 2. --- In a lattice ordered group, if x and y are coprime and if \(z \geqslant 0\) and \(x \leqslant y + z\) , then \(x \leqslant z\) .

COROLLARY 3. --- In a lattice ordered group, if x is coprime to y and to z, then it is also coprime to \(y + z\) .

COROLLARY 4. --- If \((x_i)_{1 \leq i \leq n}\) and \((y_j)_{1 \leq j \leq m}\) are two finite families of elements of a lattice ordered group G, such that each \(x_i\) is coprime to each \(y_j\) , then \(x_1 + \cdots + x_n\) is coprime to \(y_1 + \cdots + y_m\) .

This follows from Cor. 3 by induction on m and n.

COROLLARY 5. --- For any integer \(n \geqslant 0\) we have \((nx)^+ = nx^+\) and \((nx)^- = nx^-\) ; for all \(n \in \mathbf{Z}\) we have \(|nx| = |n| \cdot |x|\) .

Indeed \(nx = nx^{+} - nx^{-}\) ; since \(x^{+}\) and \(x^{-}\) are coprime, so are \(nx^{+}\) and \(nx^{-}\) if \(n \geqslant 0\) (Cor. 4); the first assertion follows by Prop. 9 b) of VI, p.~12. The second follows from the first by Prop. 9 d) in the case \(n \geqslant 0\) ; the case \(n < 0\) follows from this as a result of the relation \(|-x| = |x|\) .

PROPOSITION 11 (DIV). --- Suppose the set \(\mathcal{P}^*\) is a lattice, and let \((a_i)\) , \((b_j)\) be two finite families of elements of A. Then every gcd of \(\prod_{i} a_i\) and

\(\prod_{j} b_{j}\) divides the product \(\prod_{i,j} \gcd(a_{i}, b_{j})\) .

COROLLARY 1 (DIV). --- If a, b, c are three elements of A such that a is coprime to b, then every gcd of a and c is also a gcd of a and bc.

COROLLARY 2 (DIV) (Euclid's lemma). --- Let \(a, b, c\) be three elements of A. If \(a\) is coprime to \(b\) and divides \(bc\) , then it divides \(c\) .

COROLLARY 3 (DIV). --- If x is coprime to y and to z, then it is coprime to yz.

COROLLARY 4 (DIV). --- If \((x_{i})\) and \((y_{j})\) are two finite families of elements of A such that each \(x_{i}\) is coprime to each \(y_{j}\) , then the product of the \(x_{i}\) is coprime to the product of the \(y_{j}\) .

COROLLARY 5 (DIV). --- If \(d\) is a gcd of \(x\) and \(y\) , then \(d^n\) is a gcd of \(x^n\) and \(y^n\) for each positive integer \(n\) .

Indeed \(xd^{-1}\) and \(yd^{-1}\) are coprime (Prop. 10 (DIV)), and hence so are \(x^{n}d^{-n}\) and \(y^{n}d^{-n}\) (Cor. 4).

PROPOSITION 12. --- Let \(x_{i}\) ( \(1 \leqslant i \leqslant n\) ) be \(n\) pairwise coprime elements in a lattice ordered group. Then

\[
\sup \left(x _ {1}, \dots , x _ {n}\right) = x _ {1} + \dots + x _ {n}.
\]

This follows from the formula \(u + v = \sup(u, v) + \inf(u, v)\) (VI, p.~10, Prop. 7) by induction on n, using the fact that \(x_{i}\) is coprime to \(x_{1} + \cdots + x_{i-1}\) for \(2 \leqslant i \leqslant n\) (Cor. 4 to Prop. 11).

Remark. --- Prop. 7 of VI, p.~10 also shows that a necessary and sufficient condition for \(x\) and \(y\) to be coprime is that \(x + y = \sup(x, y)\) .

PROPOSITION 12 (DIV). --- Let \(a_i\) be a finite number \(n\) of pairwise coprime elements of \(A\) ; then the product \(a_1 \ldots a_n\) is an lcm of \(a_1, \ldots, a_n\) .

PROPOSITION 13. --- In a lattice ordered group G, let \((x_{\alpha})\) be a family having an infimum (resp. supremum), and let z be an arbitrary element of G; then the family \((\sup(z, x_{\alpha}))\) (resp. \((\inf(z, x_{\alpha})))\) has an infimum (resp. supremum) and

\[
\left\{ \begin{array}{l} \inf _ {\alpha} (\sup (z, x _ {\alpha})) = \sup \left(z, \inf _ {\alpha} x _ {\alpha}\right) \\ \sup _ {\alpha} (\inf (z, x _ {\alpha})) = \inf \left(z, \sup _ {\alpha} x _ {\alpha}\right) \end{array} \right.\tag{4}
\]

respectively.

Suppose that the family \((x_{\alpha})\) has an infimum \(y\) . We will show that \(\sup (z,y)\) is an infimum of the family \((\sup (z,x_{\alpha}))\) .

Indeed \(\sup (z, x_{\alpha}) = z + (x_{\alpha} - z)^{+}\) and by translating we can reduce to the case \(z = 0\) , that is we must show that \((x_{\alpha}^{+})\) has an infimum equal to \(y^{+}\) . Since \(y \leqslant x_{\alpha}\) we have \(y^{+} \leqslant x_{\alpha}^{+}\) for all \(\alpha\) (VI, p.~12, Prop. 9, c)). Conversely, if \(a \leqslant x_{\alpha}^{+}\) for all \(\alpha\) , then it follows that \(a \leqslant x_{\alpha} + x_{\alpha}^{-}\) (Prop. 9, a)); now \(y^{-} \geqslant x_{\alpha}^{-}\) since \(y \leqslant x_{\alpha}\) ; hence \(a \leqslant x_{\alpha} + y^{-}\) for all \(\alpha\) , that is \(a \leqslant y + y^{-} = y^{+}\) .

The other formula follows by changing to the opposite order relation.

COROLLARY. --- If an element z of a lattice ordered group G is coprime to each element \(x_{\alpha}\) of a family having an infimum y, then z is coprime to y.

This is an immediate consequence of the second formula (4).

Remark. --- Applying the formulae of Prop. 13 to a family of two elements \((x, y)\) , we obtain the following formulae, which express the fact that each of the laws of composition sup, inf in a lattice ordered group is distributive with respect to the other :

\[
\begin{array}{l} \sup (z, \inf (x, y)) = \inf (\sup (z, x), \sup (z, y)) \\ \inf (z, \sup (x, y)) = \sup (\inf (z, x), \inf (z, y)). \end{array}
\]

This distributivity property is peculiar to lattice ordered groups, and does not extend to arbitrary lattices, or even to lattice ordered monoids (cf.~VI, p.~33, Ex. 24).
% semantic-fragments: legacy-input-seam
\relax{}
